\section*{RESUMO}
O presente trabalho investiga as propriedades termodinâmicas, estruturais e cinéticas de modelos de membranas lipídicas, fundamentando aplicações na biofísica translacional. Utilizando simulações de Dinâmica Molecular com o campo de força \textit{coarse-grained} Martini 3, avaliou-se o comportamento de bicamadas compostas por DPPC e DSPC puros, em misturas binárias e em sistemas submetidos a estresses químicos por colesterol e força iônica (NaCl). A análise termodinâmica isolada revelou que a estabilidade lamelar é regida por dois eixos distintos: interações de Coulomb na interface solvatada e forças de dispersão de Lennard-Jones no empacotamento apolar. Constatou-se uma profunda assimetria cinética na transição de fase homóloga, onde o DSPC, devido à sua maior inércia estérica, permaneceu retido em um estado metaestável de líquido super-resfriado. A quantificação deste aprisionamento evidenciou a severa penalidade entálpica da barreira de dessolvatação interfacial. Adicionalmente, demonstrou-se que a marcante condensação lateral induzida pelo colesterol atua através de um mecanismo de compensação isoperibólica: o déficit eletrostático decorrente da expulsão hídrica é anulado pelo ganho massivo de contatos de van der Waals, proporcionado pelo ordenamento \textit{all-trans} das cadeias lipídicas contra a estrutura rígida do esterol. Por fim, a validação deste rigoroso balanço energético e topológico consolida a Dinâmica Molecular como uma ferramenta preditiva essencial para o desenho de vetores lipossomais e para a prospecção virológica \textit{in silico}.

\vspace{\baselineskip}
\noindent
\textbf{Palavras-chave}: Dinâmica Molecular. Biofísica Translacional. Campo de Força Martini 3. Metaestabilidade. Efeito Condensador. Termodinâmica. DPPC. DSPC.


\section*{ABSTRACT}
This work investigates the thermodynamic, structural, and kinetic properties of lipid membrane models, establishing foundations for translational biophysics. Using Molecular Dynamics simulations with the coarse-grained Martini 3 force field, the behavior of bilayers composed of pure DPPC and DSPC, binary mixtures, and systems subjected to chemical stresses by cholesterol and ionic strength (NaCl) was evaluated. Isolated thermodynamic analysis revealed that lamellar stability is governed by two distinct axes: Coulomb interactions at the solvated interface and Lennard-Jones dispersion forces in the apolar packing. A profound kinetic asymmetry was observed in the homologous phase transition, where DSPC, due to its higher steric inertia, remained trapped in a metastable supercooled liquid state. The quantification of this entrapment highlighted the severe enthalpic penalty of the interfacial desolvation barrier. Additionally, it was demonstrated that the marked lateral condensation induced by cholesterol acts through an isoperibolic compensation mechanism: the electrostatic deficit resulting from water expulsion is nullified by the massive gain in van der Waals contacts, provided by the \textit{all-trans} ordering of the lipid chains against the rigid sterol structure. Finally, the validation of this rigorous energetic and topological balance consolidates Molecular Dynamics as an essential predictive tool for the design of liposomal vectors and \textit{in silico} virological screening.

\vspace{\baselineskip}
\noindent
\textbf{Keywords}: Molecular Dynamics. Translational Biophysics. Martini 3 Force Field. Metastability. Condensing Effect. Thermodynamics. DPPC. DSPC.

